\documentclass[12pt,a4paper]{article}
\usepackage[top=2.0cm, left=2.5cm, bottom=3.0cm, right=3.0cm]{geometry}
\usepackage{times}
\usepackage[T1]{fontenc}
\usepackage[utf8]{inputenc}
\usepackage{setspace}
\usepackage{indentfirst}
\usepackage{parskip}
\usepackage{hyperref}

\setlength{\parindent}{3em}
\setlength{\parskip}{0pt}
\singlespacing

\begin{document}

\begin{center}
\textbf{\large ABSTRAK}\\[12pt]
\textbf{projectLSA: Aplikasi Shiny untuk Analisis Struktur Laten\\dengan Antarmuka Pengguna Grafis}
\end{center}

\vspace{12pt}

Invensi ini mengungkapkan suatu sistem perangkat lunak berbasis web yang diimplementasikan pada komputer untuk melakukan analisis struktur laten secara terintegrasi. Sistem ini dibangun menggunakan bahasa pemrograman R dan kerangka kerja Shiny, menyediakan antarmuka pengguna grafis interaktif yang memungkinkan pengguna melakukan berbagai metode analisis statistik laten tanpa menulis kode pemrograman. Sistem terdiri dari lima modul analitis utama: Latent Profile Analysis (LPA) untuk identifikasi subkelompok berdasarkan indikator kontinu, Latent Class Analysis (LCA) untuk klasifikasi berdasarkan indikator kategorikal, Item Response Theory (IRT) untuk pemodelan respons item, Exploratory Factor Analysis (EFA) untuk penemuan struktur faktor, dan Confirmatory Factor Analysis (CFA)/Structural Equation Modeling (SEM) untuk pengujian model pengukuran dan struktural. Sistem ini mengintegrasikan kerangka perbandingan model terstandarisasi, visualisasi interaktif, diagram jalur yang dapat dikustomisasi, dan sistem pelaporan otomatis yang menghasilkan dokumen siap publikasi. Invensi ini ditujukan untuk meningkatkan aksesibilitas, efisiensi, dan reprodusibilitas dalam analisis struktur laten di bidang penelitian pendidikan, psikometri, dan ilmu sosial.

\vspace{24pt}

\begin{center}
\textbf{\large ABSTRACT}\\[12pt]
\textbf{projectLSA: Shiny Application for Latent Structure Analysis\\with a Graphical User Interface}
\end{center}

\vspace{12pt}

The present invention discloses a web-based computer-implemented software system for conducting integrated latent structure analysis. The system is built on the R programming language and the Shiny framework, providing an interactive graphical user interface that enables users to perform multiple latent statistical analysis methods without writing programming code. The system comprises five principal analytical modules: Latent Profile Analysis (LPA) for identifying subgroups based on continuous indicators, Latent Class Analysis (LCA) for classification based on categorical indicators, Item Response Theory (IRT) for item response modeling, Exploratory Factor Analysis (EFA) for discovering factor structures, and Confirmatory Factor Analysis (CFA)/Structural Equation Modeling (SEM) for testing measurement and structural models. The system integrates a standardized model comparison framework, interactive visualization, customizable path diagrams, and an automated reporting system that generates publication-ready documents. The invention is intended to enhance accessibility, efficiency, and reproducibility in latent structure analysis for educational research, psychometrics, and social sciences.

\end{document}
